\section{Tabla de aportes}
\label{sec:aportes}

La tabla recoge los aportes publicados y las evidencias disponibles.

\begin{table}[htbp]
\centering\small
\caption{Aportes y evidencias disponibles para su revisión.}
\label{tab:aportes}
\renewcommand{\arraystretch}{1.15}
\begin{tabularx}{\linewidth}{>{\raggedright\arraybackslash}p{2.6cm}>{\raggedright\arraybackslash}X>{\raggedright\arraybackslash}p{4.8cm}}
\toprule
Integrante & Trabajo realizado & Evidencia \\
\midrule
Muñoz Vera Melanie Michelle
& Preparación del análisis estadístico y los gráficos, verificador
de Lamport, documentación del procedimiento e integración del informe
LaTeX. El uso de IA se detalla en el apartado~\ref{sec:uso-ia}.
\par Pruebas en su equipo Melanir:
capturas tomadas por Melanie, revisión de los PDF e incorporación
de resultados propios, discusión y conclusiones.
& \textbf{Publicado:}
\par 17/11: \href{https://github.com/iMouraad/PE-U1-MPZ/commit/f7cb1edbd9493d57c989a1cf0675eb7d19c0a8c7}{\texttt{f7cb1ed}}, análisis.
\par 18/11: \href{https://github.com/iMouraad/PE-U1-MPZ/commit/f6d720b9384bb439648a83e43085e53c608e8f24}{\texttt{f6d720b}}, verificador y captura;
\href{https://github.com/iMouraad/PE-U1-MPZ/commit/8b71513859cf321cc769e77babbf85aa14f0af1b}{\texttt{8b71513}}, arquitectura.
\par 19/11: \href{https://github.com/iMouraad/PE-U1-MPZ/commit/cd8f1564ffd186764369966a73be11144941897e}{\texttt{cd8f156}}, análisis de CSV reales.
\par \textbf{En este avance:} las pruebas propias del 19/11, sus seis capturas,
el programa común de medición y la integración posterior al informe.
\href{https://github.com/iMouraad/PE-U1-MPZ/commits/main/data/melanie/2026-11-19}{Historial de las pruebas}. \\
\addlinespace
Pozo Castillo Steven Fernando
& Definición del contrato Protocol Buffers y generación del código Python;
implementación de tres nodos gRPC con reloj de Lamport y registro de
eventos; medición de latencia secuencial y concurrente con solicitudes
serializadas de 100 y 10\,000 bytes. El uso de IA se detalla en el
apartado~\ref{sec:uso-ia}.
& \textbf{Publicado:}
\par 17/11: \href{https://github.com/iMouraad/PE-U1-MPZ/commit/86b5e85ca08861bdac21fcd195c7f05e2df26bcc}{\texttt{86b5e85}}, contrato y código generado.
\par 18/11: \href{https://github.com/iMouraad/PE-U1-MPZ/commit/f9d406813e207d75d3285e051fc182d5c5de641d}{\texttt{f9d4068}}, nodos gRPC, Lamport y CSV de eventos.
\par 19/11: \href{https://github.com/iMouraad/PE-U1-MPZ/commit/09ae6a6151b1b87902eeed184d9b445f4ec1622e}{\texttt{09ae6a6}}, latencia y prueba concurrente. \\
\addlinespace
Zambrano Yong Angel Daniel
& Servidor y cliente TCP con mensajes JSON delimitados por salto
de línea; programa único para los nodos A, B y C (puertos 5000 a 5002)
con un hilo servidor y envíos a destinos aleatorios; reloj de Lamport
protegido con \texttt{threading.Lock} y registro de 42 eventos en CSV.
Medición de latencia de sockets del 19/11 (1000 mediciones secuenciales
y 600 concurrentes) y sesión propia del 22/11 en el equipo Mourad con
\texttt{medir\_equipo.py}, para sockets y gRPC (4800 filas por mecanismo).
El uso de IA se detalla en el apartado~\ref{sec:uso-ia}.
& \textbf{Publicado:}
\par 17/11: \href{https://github.com/iMouraad/PE-U1-MPZ/commit/f2cda2fa390289f92c147001083e2d4a762343cc}{\texttt{f2cda2f}}, servidor y cliente TCP.
\par 18/11: \href{https://github.com/iMouraad/PE-U1-MPZ/commit/6fb25f72f4438cb471d621e66ba2cf425eb21fca}{\texttt{6fb25f7}}, nodos con Lamport; \href{https://github.com/iMouraad/PE-U1-MPZ/commit/9aff082a660cc7bc8fe8e712b92db156136886a0}{\texttt{9aff082}}, CSV de eventos.
\par 19/11: \href{https://github.com/iMouraad/PE-U1-MPZ/commit/203b1fd3ab3fcb19ec392932d67e47eec7e497e1}{\texttt{203b1fd}}, medición de latencia; \href{https://github.com/iMouraad/PE-U1-MPZ/commit/2feec51154b9d10b508b5b6a0db44be563d4bd69}{\texttt{2feec51}}, datos.
\par 22/11: \href{https://github.com/iMouraad/PE-U1-MPZ/commit/3252cc245fddf07cfb16efb3e48a4afebcca5b2a}{\texttt{3252cc2}}, mediciones propias; \href{https://github.com/iMouraad/PE-U1-MPZ/commit/482082ad641d33e25c9be1332f30d0496e3556f7}{\texttt{482082a}}, hashes.
\par Capturas 00 a 03 en \texttt{evidencias/}. \\
\bottomrule
\end{tabularx}
\end{table}

Las rutas de los datos y de sus salidas se detallan en el
apartado~\ref{sec:resultados-melanie}.
La bitácora detalla la ayuda utilizada y los problemas encontrados.
Esta tabla no atribuye a Melanie la implementación de los nodos TCP
ni de los nodos gRPC desarrollados por sus compañeros.
